\documentclass[10pt,a4paper]{amsart}
\usepackage[margin=19mm]{geometry}
\usepackage{amsmath,amssymb,mathtools}
\usepackage[scr=rsfs,cal=euler]{mathalfa}
\usepackage{xfrac}
\usepackage[unicode,hidelinks,bookmarks=false]{hyperref}
\newcommand{\ie}{i.e.\ }
\newcommand{\cf}{cf.\ }
\newcommand{\resp}{respectively\ }
\newcommand{\Subsection}{\subsection}
\providecommand{\nobreakdash}{\nobreak\hbox{-}}
\setlength{\emergencystretch}{2em}
\multlinegap0pt
\usepackage[normalem]{ulem}
\usepackage{amssymb}
\usepackage{mathtools}
\usepackage{tikz}
\newtheorem{lemma}{Lemma}
\newcommand{\LL}{{\mathchoice
{\,\LM7\,}{\,\LM7\,}{\,\LM5\,}{\,\LM{3.35}\,}}}
\newcommand{\LM}[1]{\hbox{\vrule width.2pt \vbox to#1pt{\vfill \hrule
width#1pt height.2pt}}}
\newcommand{\N}{\mathbb{N}}
\newcommand{\R}{\mathbb{R}}
\newcommand{\e}{\varepsilon}
\title{Infinitesimal Minkowskianity for manifolds with continuous Lorentzian metrics}
\author{Vanessa Ryborz}
\date{Source: arXiv:2604.22420v1 (CC BY 4.0)}
\begin{document}
\maketitle

\noindent\emph{Source and licence.} Infinitesimal Minkowskianity for manifolds with continuous Lorentzian metrics. Version arXiv:2604.22420v1 (CC BY 4.0), distributed by arXiv under the Creative Commons Attribution 4.0 International licence, http://creativecommons.org/licenses/by/4.0/, as recorded in the arXiv licence metadata for this exact version and shown on the abstract page https://arxiv.org/abs/2604.22420v1. Full licence notice: \texttt{licenses/arxiv-2604.22420-CC-BY-4.0.txt}.\par\smallskip

\noindent\textit{Editorial scope.} Counted unit: the article's lemma ball\_evaluation\_measurable (source0.tex lines 274--278). The article's complete original proof of it is delivered with it (source0.tex lines 279--350, one \texttt{proof} environment of 72 lines). No mathematics is written, repaired or re-derived: the statement and the proof are the article's own lines, in the article's own notation, and no retained line is substituted or re-flowed.  No further source-local context is needed: every cross-reference of every retained line resolves inside the two delivered spans.  Nothing was omitted from any retained span, so the package carries no omission record.  No retained line contains a citation, so the wrapper carries no bibliography and no bibliography entry.  No retained line contains a figure environment, an image-inclusion command or drawing code, so no image is delivered and the package declares no asset dependency.  Editorial formatting only, in addition: an AMS \texttt{amsart} preamble that restates the article's own declarations and macros used by the retained lines, namely the environments \texttt{lemma} on separate counters, together with the standard packages those lines need.  The retained environments print in this document as Lemma 1 in order of appearance, whereas the article prints the same items with the numbering of its own full document.  The complete native source of the article as distributed in the arXiv e-print for this exact version is bundled unchanged as \texttt{sources/199091/source0.tex}.
\begin{lemma}\label{ball_evaluation_measurable}
    Let $\Omega\subset \R^m$ be open and let $\mu$ be a locally finite Radon measure on $\Omega$. 
    For $\delta>0$ define the set $\Omega_\delta:= \{p \in \Omega: \overline{B_\delta(p)}\subset \Omega\}$. 
    Then for every $\delta>0$, the map $E_{\mu, \delta}\colon \Omega_\delta \to [0, \infty),\,   p\mapsto \mu(\overline{B_\delta(p)})$ is $\mathcal{L}^m$-measurable and locally bounded.
\end{lemma}

\begin{proof}
    Fix a $\delta>0$. 
    Fix a compact set $K_1 \subset \Omega_\delta$ and a compact set $K_2 \subset \Omega$ such that $\overline{B_\delta(K_1)} \subset K_2$. Now for each $p \in K_1$, it holds that 
    \begin{align*}
        \mu(\overline{B_\delta(p)}) \leq \mu(K_2) < \infty.
    \end{align*}
    This proves the local boundedness. 

    For $i=0, \ldots, m-1$, define the sets $P_{i}$ and the Radon measures $\mu_{i+1}, \nu_{i}$ such that 
    \begin{itemize}
        \item $P_{i}$ is a countable union of $i$-dimensional spheres, where $i$-dimensional spheres are given by $\{q \in \R^m: |q-p|=r, q-p \in V\}$, where $p \in \R^m$, $r>0$ and $V \subset \R^m$ is an $i$-dimensional linear subspace of $\R^m$. If $i=0$, $P_0$ is going to be a countable set of points.
        \item $\nu_{i}= \mu_{i}\LL P_{i}$.
        \item $\mu_{i+1} = \mu_{i}-\nu_{i}$ is zero on spheres of dimension at most $i$.
    \end{itemize}

    Define $\mu_0:= \mu$. 
    Now, define the set 
    \begin{align*}
        P_0:= \{p \in K_2: \mu(\{p\})>0\}. 
    \end{align*}
    This set has to be countable, as $P_0 = \bigcup_{n \in \N} \{p \in K_2: \mu(\{p\})>\frac{1}{n}\}$, and for each $n\in \N$, we have that $\#\{p \in K_2: \mu(\{p\})>\frac{1}{n}\}\leq n\mu(K_2) <\infty$. 
    Now define the measures
    \begin{align}
        \nu_0:= \mu \LL P_0,\ \mu_1:= \mu-\nu_0.  
    \end{align}
    Now for $i \in \{1, \ldots, m-1\}$, define inductively the measures $\nu_i, \mu_{i+1}$ as follows:
    Let
    \begin{align}
        P_i:=\{S:  i\mathrm{-dimensional\ sphere}, \ S \subset K_2,\ \mu_i(S)>0\}. 
    \end{align}
    For every $n \in \N$, we have that 
    \begin{align}\label{finitely_many_spheres_of_big_measure}
       \#\{S:  i\mathrm{-dimensional\ sphere}, \ S \subset K_2,\ \mu_i(S)>1/n\} < n\mu_{i}(K_2) \leq n \mu(K_2) < \infty.
    \end{align}
    Indeed, suppose that this was not true. Then there exist two distinct $i$-dimensional spheres $S_\alpha, S_\beta \in P_i$ such that $\mu_{i}( S_\alpha \cap S_\beta) >0$. 
    Now, we have that for any two distinct $i$-dimensional spheres, their intersection is the finite union of spheres of dimension at most $i-1$. By our assumption, the measure $\mu_{i}$ is zero on spheres of dimension at most $i-1$, which yields a contradiction and hence verifies \eqref{finitely_many_spheres_of_big_measure}.
    This proves that $P_i = \bigcup_{n \in \N} \{S:  i\mathrm{-dimensional\ sphere}, \ S \subset K_2,\ \mu_i(S)>1/n\}$ is countable. 
    Now we define 
    \begin{align}
        \nu_i:= \mu_i \LL P_i,\ \mu_{i+1}:= \mu_i-\nu_{i}. 
    \end{align}
    In total, this gives us a decomposition of the measure 
    \begin{align}
        \mu = \mu_m + \sum_{i=0}^{m-1} \nu_{i},
    \end{align}
    where the $\nu_i$ are supported on countably many $i$-dimensional spheres (collected in $P_i$) and $\mu_m$ is zero on spheres of dimension smaller than $m$. 

    For each $i\leq m-1$, and any $i$-dimensional sphere $S$, we have that the set
    \begin{align*}
        D_{S, \delta}:= \{p: S \subset \partial B_\delta(p)\}
    \end{align*}
    is a sphere of dimension $m-i-1$. Any such sphere is a Borel $\mathcal{L}^m$-zero set.
    Then, the set 
    \begin{align}
        D:= \bigcup_{i=0}^{m-1} \bigcup_{S \in P_i} D_{S, \delta}
    \end{align}
    is a Borel $\mathcal{L}^m$-zero set, which is a countable union of spheres of dimension at most $m-1$.
    For any point $p \in K_1\setminus D$, we have that 
    \begin{align*}
        \mu(\partial B_\delta(p)) = \mu_m(\partial B_\delta(p))=0.
    \end{align*}
    Note that 
    \begin{align}
        \lim_{n \to \infty} \mu(B_{\delta+ 2^{-n}}(p)\setminus B_{\delta- 2^{-n}}(p)) = \mu(\partial B_\delta(p))=0.
    \end{align}
    Fix an $\e>0$ and $n\in \N$ such that $\mu(B_{\delta+ 2^{-n}}(p)\setminus B_{\delta- 2^{-n}}(p)) \leq \e$.
    Then we have that for each $q \in B_{2^{-n-1}}(p)$, it holds that 
    \begin{align}
        |\mu(\overline{B_\delta(p)}) - \mu(\overline{B_\delta(q)})| &\leq \mu((\overline{B_\delta(p)}\setminus \overline{B_\delta(q)})\cup (\overline{B_\delta(q)}\setminus \overline{B_\delta(p)})) \leq \mu(B_{\delta+ 2^{-n}}(p) \setminus B_{\delta- 2^{-n}}(p)) \leq \e.
    \end{align}
    Hence, $E_{\mu, \delta}$ is continuous at $p$. This holds for $p \in K_1 \setminus D$, so the map is continuous $\mathcal{L}^m$-almost everywhere. Hence, $E_{\mu, \delta}$ is $\mathcal{L}^m$-measurable when restricted to $K_1$. As $K_1 \subset \Omega_\delta$ was arbitrary, this concludes the proof.
\end{proof}

\end{document}
